\section{Analysis and current status}

\begin{frame}{Planned primary analysis: MANCOVA}
\begin{columns}[T,onlytextwidth]
\column{.27\textwidth}
\textbf{Joint outcomes}
\[
\mathbf{Y}_{\mathrm{post}}=
\begin{pmatrix}
FF_{\mathrm{post}}\\ FG_{\mathrm{post}}\\ FA_{\mathrm{post}}\\ GA_{\mathrm{post}}
\end{pmatrix}
\]
\column{.69\textwidth}
\begin{itemize}
\item \textbf{Factor:} condition, with or without the additional Poincar\'e sphere.
\item \textbf{Covariates:} baseline FF, FG, FA, and GA.
\item \textbf{Target:} total condition effect on post-test outcomes (RQ1 and RQ2).
\end{itemize}
\end{columns}
\medskip
\textbf{Exploratory (RQ3):} moderation by baseline fidelities and spatial reasoning; associations with reported cognitive load.
\note{GA denotes structural/Gestalt alignment and FA functional alignment. The revised primary specification adjusts for the four baseline outcomes, not post-intervention cognitive load or spatial reasoning. Cognitive load may be affected by the intervention; ordinary adjustment would not estimate the total effect and does not automatically identify a direct effect. Spatial reasoning is measured after the applet, so its use as an aptitude moderator requires justification that treatment and differential fatigue do not affect its measurement. Baseline FG can moderate a later effect but cannot mediate it.

Keep load analyses exploratory and associational; post-session load and outcome measures alone do not identify causal mediation. Profile-specific moderation would require both baseline FF and FG, not FG alone. Reconcile the power calculation's two-covariate assumption with the revised four-baseline model. Specify follow-up directional contrasts, multiplicity, and missing-data handling before analysis. Record this revised analysis specification and any departure from registration transparently. The registration was not independently reverified and has not been modified.}
\end{frame}

\begin{frame}{Student feedback: Usefulness and usability}

\small
\begin{columns}[T,onlytextwidth]
\column{.48\textwidth}
\begin{mynote}[Perceived Usefulness]
\begin{itemize}
\item Students valued interactive graphics and seeing parameter changes immediately.
\item Hints helped, including guidance without revealing the answer.
\item Requests included more examples and fuller explanations, a glossary and definitions of polarization and circular handedness.
\end{itemize}
\end{mynote}
\column{.48\textwidth}
\begin{mynote}[Perceived Usability]
\begin{itemize}
\item Camera resets, perspective controls, and scrolling caused difficulties.
\item Requests included typed values and persistent display choices.
\item Tasks and plots were hard to view together; earlier formulas were inaccessible.
\end{itemize}
\end{mynote}
\end{columns}
\source{Summary of open-ended pilot feedback: 8 respondents across questions.}
\note{Source: pilot\_results/comments.md, read 25 September 2026. Eight distinct respondents supplied substantive comments across questions (1, 2, 3, 5, 6, 8, 9, 10); item response counts vary. These are qualitative comments, not the GETAMEL scale results. Preserve the supplied historical condition labels: without sphere = Polarization 1, with sphere = Polarization 2. The current applet uses the opposite module mapping from step 3, so its module names must not be used to relabel the pilot.

Interactive graphics and parameter manipulation appear in responses 1, 2, 6, 8, 9, and 10. Hints are mentioned positively by 3 and 9; 6 also values advice and explanations. Response 3 requests operator examples and finds explanations too short. Response 5 reports scrolling trapped over plots and formulas unreadable on smaller screens. Response 6 reports camera resets and touchpad difficulty; responses 1, 9, and 10 mention zoom or perspective controls. Response 2 requests numeric entry. Response 9 reports display selections resetting, delayed vector brackets, and task/plot separation. Responses 8 and 9 request returning to the top after advancing. Response 10 requests access to earlier formulas. These describe the pilot experience, not verified defects in the current version.}
\end{frame}

\begin{frame}{Student feedback: E-Learning self-effcacy and enjoyment}
\small
\begin{columns}[T,onlytextwidth]
\column{.48\textwidth}
\begin{mynote}[E-Learning self-efficacy]
\begin{itemize}
\item Hints enabled participants to proceed independently.
\item A student solved tasks despite uncertainty about matrices and terminology.
\item Task completion did not necessarily indicate conceptual confidence.
\end{itemize}
\end{mynote}
\column{.48\textwidth}
\begin{mynote}[Enjoyment]
\begin{itemize}
\item Students enjoyed changing the graphics and using the simulation.
\item One student compared successive tasks to a game's reward system.
\item Complex content and awkward controls reduced enjoyment.
\end{itemize}
\end{mynote}
\end{columns}
\source{Summary of Open-ended pilot feedback: 8 respondents across questions.}
\note{Response 1 describes progressing without fully understanding technical terms and uncertainty about matrices. These are two comments from one respondent, not evidence that multiple students completed tasks without understanding. Response 9 asks for a glossary, a brief polarization definition, and definitions of right/left circular polarization. Response 1 asks why the activity was done and what it has to do with quantum physics. Response 8 reports complex formulas and overly complex content. Response 6 identifies circular/rotated light as interesting; response 10 identifies the mathematical background of linear polarizers. These examples qualify any blanket statement that the mathematical material was disliked.

The reward-system analogy comes from response 1 alone. Positive comments from the sphere condition refer to interactive or three-dimensional graphics generally; none specifically identifies the sphere as the source of a benefit. This does not establish no sphere effect. Some optional field-of-study entries identify school students, so do not assume this pilot has the main study's undergraduate population. Do not infer measured self-efficacy gains, objective learning gains, or theme prevalence from these comments.}
\end{frame}

\begin{frame}{Outlook}
\textbf{Pilot study takeaways:} 
\begin{itemize}
    \item Students valued interactivity, graphics, and hints.
    \item Comments identified usability difficulties and requests for conceptual guidance.
    \item The comments do not isolate a benefit of the Poincaré sphere.
\end{itemize}
\textbf{Next steps:} 
\begin{itemize}
    \item Collect the remaining expert data and conduct the main study.
    \item As physics education experts, we need your help (open until 09th October 2026): 
\end{itemize}
\centering
\qrcode{https://nimble.li/omz88r4m} \\
\url{https://nimble.li/omz88r4m}
\end{frame}

% \begin{frame}{Initial pilot findings}
% \keyidea{Pilot studies are ongoing.}
% \begin{columns}[T,onlytextwidth]
% \column{.48\textwidth}
% \begin{mynote}[Expert pilot]
% \begin{itemize}
% \item Add sample and completion status.
% \item Add item-classification and agreement findings.
% \item Add resulting item revisions.
% \end{itemize}
% \end{mynote}
% \column{.48\textwidth}
% \begin{mynote}[Novice pilot]
% \begin{itemize}
% \item Add sample and completion status.
% \item Add interpretation or usability issues.
% \item Add resulting applet revisions.
% \end{itemize}
% \end{mynote}
% \end{columns}
% \medskip
% \placeholder{images/pilot_findings.png}{1.1cm}{Placeholder: pilot figure, applet screenshot,\\or anonymized spontaneous verbalization}
% \note{Populate this slide once pilot observations are available. Distinguish preliminary observations from established patterns. Do not present pilot findings as evidence of main-study efficacy. All bullets are placeholders.}
% \end{frame}

% \begin{frame}{Exploratory analyses: H3 and H4}
% \begin{mynote}[H3: Spatial reasoning and cognitive load]
% Regress pre/post changes in the study outcomes on spatial reasoning ability and intrinsic, extraneous, and germane cognitive load.
% \par\medskip
% \textbf{Do these characteristics help explain variation in learning?}
% \end{mynote}
% \begin{mynote}[H4: Initial Fidelity of Gestalt as a moderator]
% Examine the interaction between baseline FG and condition in predicting post-test FF or FA, controlling for the corresponding baseline outcome.
% \par\medskip
% \textbf{Does the effect of the added visualization depend on students' initial FG?}
% \end{mynote}
% \note{Retain H3 and H4 as exploratory analyses. The tracking subset and the think-aloud sample are not assumed to be identical. H3 concerns changes in the hypothesized direction: increases in FF, FA, and GA, and decreases in FG. Define change-score coding explicitly. H4 tests moderation by baseline FG. By itself, it does not identify a dual-type-specific effect, because the dual profile also requires high baseline FF. Regressions involving cognitive load provide associations and do not by themselves identify causal effects of load.}
% \end{frame}

%\begin{frame}{What would the findings tell us?}
%\begin{tabularx}{\textwidth}{@{}P{.40\textwidth}Y@{}}
%\toprule
%\textbf{Possible pattern} & \textbf{Interpretation to examine}\\
%\midrule
%Higher FA / GA; little FF / FG change & Improved expert alignment within model understanding, with stable Fidelity.\\
%Higher FF; lower FG & Consistent with a more functional stance toward models.\\
%Spontaneous tension followed by revision & Time-ordered evidence consistent with a proposed learning mechanism.\\
%Confusion or higher load without gains & Difficulties in connecting representations may limit learning.\\
%\bottomrule
%\end{tabularx}
%\source{Interpretive scenarios, not empirical results.}
%\note{A temporal sequence may support a causal explanation of an individual episode, but it does not establish population-level mediation. Confusion alone does not establish dissonance. Interpret verbalizations with actions and baseline FF/FG. Absence of verbalized conflict does not prove absence of conflict. None of these outcomes demonstrates model competency gains.}
%\end{frame}

%\begin{frame}{Contribution of the polarization study}
%\begin{itemize}
%\item Examine fidelity beliefs and expert alignment as aspects of model understanding.
%\item Test whether the additional sphere supports functional reasoning and changes visual-fidelity beliefs.
%\item Use spontaneous think-aloud observations to investigate learning and difficulty during existing tasks.
%\end{itemize}
%\begin{myalert}[Scope of the conclusions]
%No additional interview tasks or follow-up questions. No direct measurement of model competencies.
%\end{myalert}
%\keyidea{Can an additional representation change both what learners understand and what they think a model is for?}
%\end{frame}
